\section{Conclusion}

The \textbf{Mario Game Project} successfully demonstrates the end-to-end architectural design, mathematical modeling, and software engineering of a fully featured 2D platformer engine developed in C++ using the SFML multimedia framework. Inspired by the foundational mechanics of classic side-scrolling platformers, the system integrates complex gameplay systems—including multi-hero selection (Mario, Luigi, and Flash), multi-tier character evolution (Small, Giant, and Tier-3 Special forms), dynamic AABB collision kinematics, custom artificial intelligence, Tiled map level parsing, spatial partitioning, interactive world elements, HUD metrics, contextual audio streams, and JSON-driven game state persistence.

\subsection{Architectural Achievements and Design Integrity}
A core objective of this project was applying Object-Oriented Programming (OOP) principles and architectural design patterns to resolve the inherent complexity of real-time interactive software. The application of these methodologies yielded substantial benefits across the codebase:
\begin{itemize}
    \item \textbf{Encapsulation and Separation of Concerns:} By strictly decoupling the core render loop (\texttt{GameLoop}), physics mechanics (\texttt{WorldPhysicsSystem}), presentation components (\texttt{HUD}, \texttt{SoundManager}), and world state management (\texttt{LevelRuntime}), individual subsystems were built and maintained independently without propagating collateral bug regressions.
    \item \textbf{Polymorphism and Behavioral Patterns:} The \textit{State/Strategy Pattern} was successfully leveraged to govern hero state transitions and dynamic texture routing without violating the Open/Closed Principle (OCP). The \textit{Factory Pattern} abstracted entity instantiation during map loading, allowing new blocks, items, and character classes to be introduced seamlessly. The \textit{Observer Pattern} cleanly decoupled engine event processing from real-time HUD rendering updates, while the \textit{Template Method Pattern} established reusable patrol and decision-making skeletons for enemy dynamic behavior hierarchies.
\end{itemize}

\subsection{Engineering Bottlenecks and Solutions}
Throughout the development lifecycle, overcoming real-world technical bottlenecks provided deep practical insight into graphics rendering, continuous collision detection, numerical mechanics, and persistent serialization:
\begin{itemize}
    \item \textbf{Mathematical Origin Alignment:} Formulated an automated pixel-alpha Center-of-Mass (CM) padding algorithm to eliminate horizontal sprite jittering and state transition sliding across non-uniform character asset sheets.
    \item \textbf{Kinematic Decoupling in FSM:} Redesigned locomotion state guards by constructing composite input-velocity invariants, permanently resolving high-frequency Finite State Machine (FSM) state oscillations and frame-freezing defects.
    \item \textbf{Map-Agnostic Dynamic Physics:} Replaced fragile hardcoded ground assumptions with continuous dynamic gravity acceleration and dual-pass AABB spatial snapping within \texttt{World\-Physics\-System}, ensuring robust interaction across arbitrary topographies.
    \item \textbf{State Serialization Pipeline:} Architected a deterministic base-reload and JSON overlay mechanism in \texttt{SaveManager}, successfully preserving dynamic entity states, damaged block topologies, player progress metrics, and invulnerability timers.
    \item \textbf{Dynamic Spatial Partitioning:} Implemented context-aware \texttt{PlayableRegion} bounding zones that dynamically clamp camera viewports, eliminate visual border bleeding, adjust underground kill-planes, and switch environmental theme palettes and audio tracks seamlessly.
\end{itemize}

\subsection{Team Synergy and Technical Growth}
The successful completion of the project is attributed to structured collaboration and a clear division of responsibilities among all four team members—spanning Core Engine \& Physics (Tran Minh Khoa), Hero \& Item Mechanics (Tran Nhu Khai), Enemy AI Subsystems (Do Viet Hoang Long), and Level Parsing, UI/UX \& Audio (Tran Dang Khoa). Cooperative version control via Git, structured peer code reviews, and continuous integration testing ensured overall system stability and code consistency throughout the development cycle.

In summary, the Mario Game Project stands as a robust, highly extensible, and production-ready C++ software system. Beyond delivering a smooth and interactive gameplay experience, the project successfully bridged theoretical computer science principles—such as OOP design patterns, state machines, numerical physics simulation, and spatial partitioning—with practical, real-world software architecture.